\documentclass[12pt]{article}

\pagestyle{empty}
\setlength{\topmargin}{0in}
\setlength{\headheight}{0in}
\setlength{\topsep}{0in}
\setlength{\textheight}{9in}
\setlength{\oddsidemargin}{0in}
\setlength{\evensidemargin}{0in}
\setlength{\textwidth}{6.5in}

\usepackage{palatino, graphicx}

\newcommand{\real}{{\bf R}}
\newcommand{\vs}[1]{\vspace{#1in}}

\begin{document}

\noindent
{\bf Mathematics 202} \\ 
{\bf Working with mass}
\bigskip

\begin{enumerate}
\item Suppose that we have a thin beam laid out horizontally on the
  $x$-axis, like the one shown below, and that the density is
  constantly $\rho = 3$ grams per 
  centimeter and that the beam is 10 centimeters long.  What is the
  mass of the beam?  What is the general recipe for finding the mass
  if you know the length and the density?

  \includegraphics{figures/19-beam.eps}

  \vs{0.5}
  Suppose now that the beam is composed of a material whose density 
  varies from left to right.  This means that the density depends on
  $x$, and we have $\rho(x)=x+3$ where the beam is between $0\leq
  x\leq 10$.  If we take a slice of the beam at position $x$ and width
  $dx$, as shown below, what is the amount of mass $dm$ in the slice?

  \includegraphics{figures/19-beamslice.eps}

  \vs{0.5}
  Now that we know the mass of each slice, how may we find the total
  mass?  What is the total mass of this beam?

  \vs{1}
\item Suppose that our beam has density $\rho(x) = e^{-x/10}$ where
  $0\leq x \leq 10$.  What is the mass of the beam?

  \vs{1}
  \newpage
\item Let's now look at a two-dimensional example and suppose that we
  have manufactured a triangular plate shown below on the left.  The
  density of the material varies from left to right as $\rho(x) = x$
  grams per square centimeter.
  \begin{center}
    \includegraphics{figures/19-plate.eps} \hspace*{0.5in}
    \includegraphics{figures/19-plateslice.eps} \\
  \end{center}
  Suppose we take a vertical slice of width $dx$ as shown on the
  right.  What is the area $dA$ of this slice?  What is its mass?

  \vs{1}
  What is the total mass of the plate?

  \vs{1}
\item Suppose we make a bowl by spinning the portion of the graph of
  the function $y=x^2/4$, where $0\leq x \leq 4$ about the $y$-axis.
  Find the volume $dV$ of a horizontal slice of thickiness $dy$ and use
  this to find the total volume of the bowl.

  \includegraphics{figures/19-solid.eps}

  \newpage
  Suppose that we fill the bowl with a liquid
  and that some sediments in the liquid settle out over time so that
  the density is given by $\rho(y) = 8-y$ grams per unit volume.
  Since you know $dV$, the volume of a horiztonal slice, what
  is the mass of a horizontal slice of the liquid?

  \vs{1}
  What is the total mass of the liquid in the bowl?

  \vs{1}
  What is the average density of the liquid?

  \vs{1}
  Now suppose that we scoop off the top of the liquid removing half of
  its volume.  What is the depth of the liquid now?

  \vs{2}
  What is the mass of the remaining liquid?

\end{enumerate}
\end{document}

  
\item Suppose that we construct a solid by spinning the portion of the
  graph of $y=1/\sqrt{x}$ between $1\leq x \leq L$ about the
  $x$-axis.  What is the volume of a vertical slice of thickness $dx$?

  \vs{1}
  What is the total volume of the solid?

  \vs{1}
  What happens to the volume as $L\to\infty$?

  \vs{1}
  Suppose that the density of the solid is $\rho(x) = 1/x$ grams per
  cubic centimeter.  What is the mass of a vertical slice of thickness
  $dx$?

  \vs{1}
  What is the mass of the solid?

  \vs{1}
  What happens to the mass as $L\to\infty$?

  \vs{1}
  After we have let $L\to\infty$, suppose we want to cut the solid
  into two pieces, each of which have the same amount of mass.  Where
  should we cut the solid?

  

\end{enumerate}

\end{document}



